\documentclass[11pt,a4paper]{article}

\usepackage[spanish,es-tabla]{babel}
\usepackage[utf8]{inputenc}
\usepackage[T1]{fontenc}
\usepackage{geometry}
\usepackage{amsmath}
\usepackage{booktabs}
\usepackage{array}
\usepackage{enumitem}
\usepackage{hyperref}
\usepackage{xcolor}
\usepackage{titlesec}
\usepackage{float}

\geometry{margin=2.5cm}
\hypersetup{
    colorlinks=true,
    linkcolor=black,
    urlcolor=blue,
    citecolor=black
}
\setlist[itemize]{noitemsep, topsep=2pt}
\emergencystretch=3em
\titleformat{\section}{\large\bfseries}{\thesection.}{0.5em}{}
\titleformat{\subsection}{\normalsize\bfseries}{\thesubsection.}{0.5em}{}

% Entorno para las conclusiones clave, imitando los bloques ```text``` de los
% correos a la tutora y de la documentación interna del proyecto.
\usepackage{tcolorbox}
\newtcolorbox{conclusion}{colback=gray!6, colframe=gray!40, boxrule=0.4pt,
  left=6pt, right=6pt, top=4pt, bottom=4pt, fontupper=\footnotesize}

\title{\textbf{Deep Learning Aplicado a la Microestructura Financiera:\\Predicción a Corto Plazo en Polymarket}}
\author{Marc Maldonado Lorca\\Universidad Politécnica de Madrid --- Máster en Deep Learning\\Trabajo Fin de Máster}
\date{Curso 2025/2026}

\begin{document}

\maketitle

\begin{abstract}
Este Trabajo Fin de Máster estudia si las técnicas de Deep Learning pueden
aprovechar señales de microestructura en los mercados de Bitcoin de Polymarket.
No se trata de predecir el resultado final de un evento, sino movimientos
locales de precio del contrato en ventanas de segundos, usando el libro de
órdenes y referencias externas de BTC. La aportación principal no es un modelo
ganador, sino una \emph{metodología} completa de captura de datos, validación
temporal y control de costes, y la honestidad con la que se aplica. La lección
central es que un baseline que acierta la dirección en torno al 90\,\% del tiempo
deja de ser rentable cuando se introduce una latencia de entrada realista: el
coste y la ejecución, no el acierto direccional, deciden el resultado. A partir
de ahí se recorre una escalera de modelos (baseline tabular, modelos
latency-aware, secuencias y orderbook con redes profundas) hasta un especialista
tabular congelado con un gate de volatilidad, que sobrevive a cinco días de
validación fuera de muestra (\,$+1{,}069$ ticks netos por operación, 5/5 días
positivos\,). Todos los resultados se expresan en ticks netos, nunca en dólares,
y se acompañan siempre de su soporte e incertidumbre. No se afirma rentabilidad
real: se entrega la metodología, la evidencia y sus límites.
\end{abstract}

\tableofcontents
\newpage

% =====================================================================
\section{INTRODUCCIÓN Y MOTIVACIÓN}

\subsection{Planteamiento}

Este trabajo continúa la propuesta aceptada en la asignatura de Proyectos de
Deep Learning. La idea de partida es tratar los mercados de Polymarket como un
pequeño exchange financiero: hay libro de órdenes, bids, asks, spread y
operaciones que llegan en el tiempo. La diferencia respecto a un exchange
tradicional es que el activo es un contrato binario sobre un evento. Me centro
en mercados BTC porque permiten cruzar la información de Polymarket con
referencias externas de Bitcoin (Binance spot y perpetuo, Chainlink).

La pregunta de fondo no es si el evento resolverá con \textsc{yes} o \textsc{no},
sino si el precio del contrato contiene señales explotables a muy corto plazo.
Para cada instante $t$ se construye una ventana con el estado del mercado y se
intenta anticipar el movimiento en un horizonte $H$. Desde el principio asumo
que una métrica de clasificación puede engañar: un modelo puede tener buen
acierto y aun así no servir si opera tarde, si aprende artefactos del dataset o
si ignora el coste real de la orden.

\subsection{Objetivo y contribución}

El objetivo declarado en la propuesta era construir un flujo reproducible,
definir controles de calidad de datos, entrenar baselines, evaluar Deep Learning
sobre secuencias de orderbook y medir el rendimiento con métricas económicas, no
solo de clasificación. Cumplidos esos pasos, la contribución defendible de la
memoria es doble:

\begin{itemize}
  \item un \textbf{sistema de datos propio} multi-fuente con gates de calidad por
  sesión, y
  \item un \textbf{protocolo de evaluación honesto} que muestra dónde y por qué se
  rompe la cadena entre predicción y ejecución, y que termina en un candidato
  congelado validado fuera de muestra bajo costes.
\end{itemize}

Es importante el matiz: la honestidad del protocolo (congelar antes de mirar,
distinguir simulación de validación, admitir los errores) es la tesis, por
encima de cualquier número concreto.

\subsection{Cómo leer esta memoria}

La memoria sigue el patrón de las prácticas del máster: baseline sencillo
$\rightarrow$ modelo con features $\rightarrow$ modelo profundo $\rightarrow$
modelo final, justificando cada salto con evidencia. El capítulo de experimentos
recorre esa escalera en orden cronológico, tal y como ocurrió, incluyendo los
caminos que no funcionaron, porque los negativos son parte del resultado.

% =====================================================================
\section{TRABAJO RELACIONADO}

El trabajo se apoya en cuatro áreas. Los mercados de predicción se estudian como
mecanismos de agregación de información: Wolfers y Zitzewitz
\cite{wolfers2004prediction} muestran que pueden producir previsiones precisas,
aunque su interpretación depende del diseño del contrato, la liquidez y los
incentivos. Esto encaja con Polymarket, donde el precio de un contrato binario
se lee como una probabilidad implícita \cite{polymarket_prices}, pero su
documentación lo describe además como un CLOB con matching off-chain y
liquidación on-chain \cite{polymarket_trading}, lo que permite tratarlo como un
exchange con spread, profundidad y slippage.

Sobre libros de órdenes, DeepLOB \cite{zhang2019deeplob} propone convoluciones
para la estructura espacial del libro y capas recurrentes para la dinámica
temporal; Sirignano y Cont \cite{sirignano2019universal} argumentan que existen
características universales en la formación de precios aprendibles con Deep
Learning, y Jha et al. \cite{jha2020digital} aplican redes convolucionales
temporales a orderbooks de criptoactivos. Como problema de series temporales,
arquitecturas como el Temporal Fusion Transformer \cite{lim2021tft} sirven de
referencia para combinar covariables e historiales.

Finalmente, la evaluación realista es un eje del trabajo. La analogía con
Pendlebury et al. \cite{pendlebury2019tesseract} sobre sesgos espaciales y
temporales es directa: si se mezclan datos futuros con el entrenamiento, el
resultado parece mejor de lo que será en producción. Las recomendaciones de Arp
et al. \cite{arp2022dos} sobre errores comunes en ML aplicado, y la deuda técnica
descrita por Sculley et al. \cite{sculley2015debt}, guían las decisiones de
protocolo y trazabilidad.

% =====================================================================
\section{SISTEMA DE DATOS}

\subsection{Recolector multi-fuente}

Construí un recolector propio que captura, por sesión de mercado, el orderbook
de Polymarket (CLOB), trades, metadata de mercado, y referencias externas de
Binance spot y perpetuo y de Chainlink. Las fuentes se alinean temporalmente
sobre una malla de dos segundos. La elección de la malla no es cosmética:
condiciona el contrato de latencia que se usa más adelante, porque en vez de
interpolar un instante que no observamos, evaluamos contra la siguiente
fotografía disponible.

\subsection{Calidad por sesión}

Cada sesión pasa controles de calidad antes de entrar al corpus: cobertura,
continuidad, ratio de datos \emph{stale}, validación de arranque (\emph{smoke})
y huecos severos. Las sesiones que no superan los umbrales se rechazan. Esta
separación entre sesiones válidas y rechazadas es deliberada: un dato externo
retrasado o incompleto puede generar una señal falsa, y prefiero perder muestras
a contaminar el entrenamiento.

\subsection{Corpus}

El corpus principal cubre del 11 al 25 de mayo de 2026 (entrenamiento, validación
y test temporal), con del orden de 2,6 millones de filas cross-venue. Como
conjunto fresco \emph{fuera de muestra} se reservó del 6 al 10 de junio de 2026,
capturado después de congelar el candidato (véase la Sección~\ref{sec:metodo}).
La unidad de trabajo es \texttt{session\_id + token\_id + time\_index\_ns} sobre
la malla de dos segundos; el esquema de la consulta base se publica en el
repositorio (\texttt{sql/core\_training\_base.sql}). El corpus completo y la
fuente SQLite del recolector no se publican por tamaño y privacidad.

% =====================================================================
\section{METODOLOGÍA}
\label{sec:metodo}

\subsection{Target y unidad de acción}

La etiqueta direccional se calcula \emph{offline} a partir del movimiento futuro
del mid-price,
\[
\Delta_H = \frac{mid_{t+H} - mid_t}{mid_t},
\]
con un umbral $\theta$ que debe cubrir spread, comisiones y ruido. Las
observaciones son causales; las etiquetas se generan a posteriori y con
trazabilidad, para no contaminar las features con información del futuro. Todo el
resultado económico se expresa en \textbf{ticks netos}, la unidad natural del
libro, y nunca en dólares: hablar en dólares introduce supuestos de tamaño de
posición y suena a sobreafirmación.

\subsection{Validación temporal y fuga de datos}

La partición es estrictamente temporal: se entrena con sesiones antiguas, se
valida con sesiones posteriores y se abre el test terminal al final. En los
experimentos secuenciales se usa además un protocolo \emph{out-of-fold} (OOF):
cada modelo se entrena solo con días anteriores, las predicciones que consume una
segunda etapa son OOF, y la política se elige sin mirar el test. Esta disciplina
contra la fuga temporal es la que evita que un resultado offline brillante se
desmorone en producción \cite{pendlebury2019tesseract}.

\subsection{Modelo de costes y latencia}

Medí la latencia real del flujo de orden de Polymarket en una prueba controlada:
la mediana de tiempo desde que se observa la oportunidad hasta el ack/match del
CLOB fue de aproximadamente $0{,}43$ segundos en p95 para la muestra válida. Por
prudencia, el contrato de modelado fija:

\begin{conclusion}
\begin{verbatim}
latencia conceptual de modelado = 1 segundo
latencia observable en el dataset = 2 segundos
\end{verbatim}
\end{conclusion}

La latencia observable de dos segundos es conservadora porque coincide con la
malla del dataset. El coste de referencia es de $0{,}5$ ticks por operación, con
un contraste optimista de $0{,}25$.

\subsection{Protocolo de congelación}

El candidato final se \textbf{congeló} (modelo, features y umbrales de política)
\emph{antes} de ver los datos fresh del 6--10 de junio. Esto convierte la
validación fresh en una prueba fuera de muestra genuina y no en otra ronda de
ajuste. Es una decisión metodológica deliberada: una vez fijada la política, no
se retoca para mejorar los números.

% =====================================================================
\section{EXPERIMENTOS: LA ESCALERA DE MODELOS}

\subsection{Baseline tabular H16}

El primer baseline tabular sobre el horizonte corto acierta la dirección en torno
al 90\,\% y pasa los folds walk-forward. Visto solo con métricas de clasificación,
parece excelente. Este es exactamente el espejismo que la metodología busca
evitar.

\subsection{El stress de latencia (lección central)}

Al introducir la entrada retrasada y el coste, el panorama cambia. La
Tabla~\ref{tab:stress} muestra el test terminal (el bloque que no se usa para
elegir la política) bajo distintas latencias, en el escenario conservador con
filtro de spread.

\begin{table}[H]
\centering
\small
\begin{tabular}{rrrrrr}
\toprule
\textbf{Latencia} & \textbf{Acciones} & \textbf{Hit up} & \textbf{Wrong down} & \textbf{Neto medio} & \textbf{Neto total} \\
\midrule
$0$\,s & 64 & 92,19\,\% & 1,56\,\%  & $+6{,}79$  & $+434{,}44$ \\
$2$\,s & 63 & 39,68\,\% & 44,44\,\% & $-1{,}48$  & $-93{,}48$  \\
$4$\,s & 66 & 45,45\,\% & 43,94\,\% & $-0{,}74$  & $-48{,}83$  \\
$8$\,s & 66 & 43,94\,\% & 31,82\,\% & $-0{,}30$  & $-19{,}53$  \\
\bottomrule
\end{tabular}
\caption{Test terminal bajo latencia (ticks). La señal a $0$\,s es real como
markout, pero no sobrevive como política ejecutable a $2$\,s.}
\label{tab:stress}
\end{table}

\begin{conclusion}
\begin{verbatim}
Acertar la dirección no es suficiente; con latencia realista de 2 s el neto
del test terminal cae a -1,48 ticks. El problema central pasa a ser la
ejecución, no la predicción.
\end{verbatim}
\end{conclusion}

Esta tabla reorganiza todo el proyecto. A partir de aquí, cualquier modelo se
juzga por su política económica bajo coste y latencia, no por su acierto.

\subsection{Modelos latency-aware}

El siguiente peldaño son modelos que incorporan la latencia en el etiquetado y la
selección. Un modelo latency-aware llegó a un AUC de test de $0{,}79$, pero la
política económica seguía siendo negativa. La conclusión es coherente con la
anterior: más capacidad de clasificación no implica una política rentable.

\subsection{Secuencias y orderbook con Deep Learning}

Construí secuencias de hasta ocho snapshots y un tensor del orderbook top-10 de
forma $5\,831 \times 8 \times 10 \times 10$, auditado como completo y alineado.
Sobre esa representación probé modelos tabulares secuenciales, GRU, una Conv1D
solo con orderbook (siguiendo la intuición de DeepLOB \cite{zhang2019deeplob}),
fusiones de secuencia tabular y orderbook, y selectores de riesgo. La GRU y la
fusión mostraron señal de representación, y algunas políticas proxy fueron
positivas, pero dejaron de serlo al exigir selección causal, estabilidad diaria o
fills parciales/adversos. Por eso no se consideran modelos finales.

\begin{conclusion}
\begin{verbatim}
El Deep Learning aporta representación, pero no una política rentable y estable
con ~15 días de datos. Es un negativo honesto, no un fracaso oculto.
\end{verbatim}
\end{conclusion}

\subsection{Corrección metodológica del score adverse}

Durante la auditoría de selección detecté un error semántico: el target
\texttt{adverse=1} representa un resultado malo, pero se había usado su
probabilidad como score descendente. La dirección correcta es
\[
\texttt{safe\_adverse\_score} = 1 - P(\text{adverse}).
\]
Marqué los resultados afectados como superados y repetí la validación con la
dirección corregida; la candidata anterior dejó de ser evidencia válida. Incluyo
este episodio deliberadamente: muestra por qué no basta con una métrica positiva,
también hay que auditar el significado económico del score. Es, probablemente, la
mejor evidencia de rigor del proyecto.

\subsection{Protocolo OOF estricto y selección de horizonte}

Repetí la fusión bajo un protocolo estrictamente out-of-fold. La segunda etapa no
produjo una política robusta (\texttt{NO\_GO\_CLEAN\_OOF\_FUSION\_STAGE2}), pero
el modelo H60 conservó una señal modesta (AUC de test $0{,}5656$; top 35\,\% tras
coste $0{,}5$ de $+0{,}1778$ ticks proxy). Para comprobar si el problema era
predecir demasiado pronto, comparé H60 frente a H120 con exactamente las mismas
secuencias (Tabla~\ref{tab:horizon}).

\begin{table}[H]
\centering
\small
\begin{tabular}{lrr}
\toprule
\textbf{Métrica} & \textbf{H60 emparejado} & \textbf{H120} \\
\midrule
AUC test                  & $0{,}5651$ & $0{,}4951$ \\
Top 35\,\% test, coste $0{,}5$ & $+0{,}4262$ & $-1{,}1171$ \\
Top 35\,\% test, coste $1{,}0$ & $+0{,}0579$ & $-1{,}4910$ \\
\bottomrule
\end{tabular}
\caption{Selección de horizonte (ticks). H120 funciona prácticamente como azar;
H60 es el único que conserva señal. H240 colapsa todavía más en fresh.}
\label{tab:horizon}
\end{table}

H120 (y H240) quedan descartados con evidencia: alargar el horizonte no mejora la
señal. El multi-horizonte se cierra y H60 se mantiene como único candidato.

\subsection{Especialista prestart H60 (candidato final)}

El candidato final es un especialista para la ventana \emph{prestart}: un
HistGradientBoosting que combina un regresor de valor esperado (EV) y un
clasificador de salud de sesión, sobre el conjunto de 36 features
\texttt{no\_clock} (sin variables de reloj, para evitar fuga temporal trivial).
La señal se restringe al filtro \texttt{strict\_45\_60\_early}: solo la ventana
de $-60$ a $-45$ segundos antes de la apertura, con libro sano y coste visible
acotado. La política congelada es:

\begin{conclusion}
\begin{verbatim}
ev_pred > 0.75   (regresor de valor esperado)
hp_pred >= 0.50  (clasificador de sesión "healthy")
filtro strict_45_60_early  (ventana prestart -60..-45 s)
\end{verbatim}
\end{conclusion}

\subsection{Vol gate}

Al preparar la validación fresh observé un cambio de régimen: la fracción de
sesiones de alta volatilidad pasó del 18--22\,\% en el corpus de entrenamiento a
cerca del 58\,\% en fresh (Tabla~\ref{tab:regime}). El modelo se entrenó en un
régimen de baja volatilidad y sus predicciones en alta volatilidad no son
fiables en la distribución desplazada. La respuesta es un gate simple de
volatilidad: operar solo cuando \texttt{perp\_realized\_vol\_bps\_5s} $\le 0{,}6657$
(el cuantil 80 de la volatilidad de entrenamiento).

\begin{table}[H]
\centering
\small
\begin{tabular}{lrl}
\toprule
\textbf{Periodo} & \textbf{Fracción HIGH\_VOL} & \textbf{Efecto del gate} \\
\midrule
train (11--20 may) & 22\,\% & mínimo (el gate casi no actúa) \\
val (21--22 may)   & 21\,\% & mínimo \\
test (23--25 may)  & 18\,\% & mínimo \\
fresh (6--10 jun)  & 58\,\% & \textbf{crítico ($+0{,}73$ de mejora)} \\
\bottomrule
\end{tabular}
\caption{El cambio de régimen entre entrenamiento y fresh justifica el vol gate.}
\label{tab:regime}
\end{table}

% =====================================================================
\section{RESULTADOS}

\subsection{Validación fresh out-of-sample}

Sobre los cinco días fresh (6--10 de junio), con el candidato congelado y el vol
gate, la validación fuera de muestra da:

\begin{conclusion}
\begin{verbatim}
n = 318 trades, 5 días (6-10 jun), ~63,6 trades/día
Net@0.5  = +1,069 ticks/trade      Net@0.25 = +0,819 ticks/trade
IC 90%   = [+0,472, +1,654]        P(neto > 0) = 99,8%
Días positivos = 5/5               Drawdown máx = 88,1 ticks
Sin vol gate: +0,349  ->  el gate aporta +0,72 ticks
\end{verbatim}
\end{conclusion}

El intervalo de confianza es bootstrap (5\,000 remuestreos, semilla fija) y la
probabilidad de neto positivo se estima sobre esa misma distribución. El soporte
son cinco días: es corto y se indica siempre junto al resultado. El gate aporta
la mayor parte de la mejora precisamente porque el régimen fresh cambió; sin él,
el neto cae a $+0{,}349$.

\subsection{Negativos valiosos: adaptación al régimen}

Probé varias formas de adaptar el modelo al nuevo régimen: reentrenamiento
walk-forward continuo, CORAL, gate conformal, ensemble bootstrap, vol gate
adaptativo, features normalizadas por volatilidad y distillation de un TCN. Todas
resultaron \texttt{NO\_GO}: ninguna superó de forma robusta al modelo congelado
con el gate simple. La conclusión es interesante por sí misma:

\begin{conclusion}
\begin{verbatim}
Con pocos días de datos, un modelo congelado + un gate simple supera a la
adaptación sofisticada. Adaptarse a un régimen exige más datos de los que hay.
\end{verbatim}
\end{conclusion}

\subsection{Paper shadow y gates de decisión}

La infraestructura final es un \emph{paper shadow ledger}: registro operación a
operación, resumen diario y gates de promoción cuantitativos. Los gates son
deliberadamente conservadores y exigen soporte, no solo neto positivo:

\begin{itemize}
  \item \texttt{PAPER\_SHADOW\_CANDIDATE}: net@0.5 $> 0{,}5$ y $\ge 200$ trades y $\ge 8$ días.
  \item \texttt{BOT\_CANDIDATE}: net@0.5 $> 0{,}5$ y $\ge 400$ trades y $\ge 12$ días.
\end{itemize}

Con cinco días, el candidato queda en \texttt{PROMISING\_NEEDS\_DATA}: por falta
de soporte temporal, no por falta de señal. Esto es exactamente lo que la
metodología debe reportar; forzar una promoción aquí sería deshonesto.

% =====================================================================
\section{CONSIDERACIONES ÉTICAS}

Este sistema toma decisiones económicas automatizadas y es especialmente
sensible al sobreajuste cuando hay pocos datos, dos riesgos que ya señalé en el
trabajo de la asignatura de Ética sobre este mismo proyecto. La memoria es
coherente con aquello: el proyecto se detiene deliberadamente en paper shadow y
no opera con capital real.

Sigo el marco de confidencialidad, integridad y disponibilidad. La integridad de
los datos me preocupa tanto como el modelo: un dato externo retrasado o
incorrecto puede generar una señal falsa y, con ella, una pérdida. Por eso el
sistema entrena solo con sesiones aceptadas, separa datos crudos de features y
labels, mantiene trazabilidad de artefactos y define límites para no operar con
datos \emph{stale}. Las recomendaciones de Arp et al. \cite{arp2022dos} ---definir
bien el problema, cuidar el etiquetado, usar baselines, evitar comparaciones
injustas--- se aplican a lo largo de todo el trabajo. El proyecto no usa datos
personales: solo datos públicos de mercado.

% =====================================================================
\section{CONCLUSIONES, LÍMITES Y TRABAJO FUTURO}

\subsection{Conclusión}

\begin{conclusion}
\begin{verbatim}
Se construyó una metodología completa de captura, validación temporal y control
de costes para mercados BTC de Polymarket. Los baselines direccionales mueren
con latencia realista; el deep learning aporta representación pero no una
política rentable con ~15 días de datos. Un especialista tabular congelado con
gate de volatilidad sí sobrevive a 5 días out-of-sample puros (+1,07 ticks
netos/trade, 5/5 días positivos) y queda en seguimiento paper-shadow con gates
cuantitativos hacia un posible bot. No se afirma rentabilidad real: se entrega
la metodología, la evidencia y sus límites.
\end{verbatim}
\end{conclusion}

\subsection{Límites}

El límite más importante es el soporte temporal: cinco días fresh. El intervalo
de confianza es estrecho y los cinco días son positivos, pero cinco días no son
una validación definitiva, y así se reporta. Además, el resultado es específico
de la ventana prestart y del régimen capturado; un cambio de régimen mayor
podría invalidar el gate, como el propio cambio entre entrenamiento y fresh
demuestra.

\subsection{Trabajo futuro}

Las líneas abiertas son: ampliar el soporte fresh hasta alcanzar los gates de
\texttt{PAPER\_SHADOW\_CANDIDATE}; validar un patrón en U observado en la
volatilidad dentro del gate (la volatilidad muy baja concentra la mayor parte del
PnL), que de momento se deja como \emph{hipótesis} y no como resultado por
descansar en solo cinco días; y revisitar el Deep Learning secuencial con 60 o
más días de datos, condicionado a que supere al especialista tabular bajo el
mismo protocolo.

% =====================================================================
\section{REPRODUCIBILIDAD}

El repositorio asociado replica la estructura de un proyecto de MLOps visto en el
máster: un paquete \texttt{src/edgerunner/} con la política congelada, el modelo
de costes y los gates extraídos y testeados (pytest); 18 documentos canónicos del
arco; 8 notebooks pedagógicos ejecutados; los scripts reproducibles de cada
peldaño; el esquema SQL; y los registros de resultados (\texttt{key\_results.csv}
y \texttt{decision\_register.csv}) filtrados al arco. Los umbrales congelados
viven en \texttt{config/frozen\_policy\_thresholds.yaml} y no se modifican en
código. El registro de experimentos y decisiones reduce la deuda técnica
inherente a estos sistemas \cite{sculley2015debt}: una pequeña diferencia en
sampling o en labels puede cambiar por completo el resultado, y por eso queda
trazada.

% =====================================================================
\begin{thebibliography}{99}

\bibitem{polymarket_prices}
Polymarket Documentation. \textit{Prices \& Orderbooks.}
\url{https://docs.polymarket.com/polymarket-learn/trading/using-the-orderbook}

\bibitem{polymarket_trading}
Polymarket Documentation. \textit{Order Lifecycle.}
\url{https://docs.polymarket.com/concepts/order-lifecycle}

\bibitem{wolfers2004prediction}
J. Wolfers and E. Zitzewitz. \textit{Prediction Markets.}
Journal of Economic Perspectives, 18(2), 107--126, 2004.
\url{https://www.nber.org/papers/w10504}

\bibitem{zhang2019deeplob}
Z. Zhang, S. Zohren and S. Roberts.
\textit{DeepLOB: Deep Convolutional Neural Networks for Limit Order Books.}
IEEE Transactions on Signal Processing, 67(11), 3001--3012, 2019.
\url{https://arxiv.org/abs/1808.03668}

\bibitem{sirignano2019universal}
J. Sirignano and R. Cont.
\textit{Universal features of price formation in financial markets: perspectives from deep learning.}
Quantitative Finance, 19(9), 1449--1459, 2019.

\bibitem{jha2020digital}
R. Jha et al. \textit{Deep Learning for Digital Asset Limit Order Books.}
arXiv:2010.01241, 2020.

\bibitem{lim2021tft}
B. Lim, S. O. Arik, N. Loeff and T. Pfister.
\textit{Temporal Fusion Transformers for Interpretable Multi-horizon Time Series Forecasting.}
International Journal of Forecasting, 37(4), 1748--1764, 2021.
\url{https://arxiv.org/abs/1912.09363}

\bibitem{pendlebury2019tesseract}
F. Pendlebury, F. Pierazzi, R. Jordaney, J. Kinder and L. Cavallaro.
\textit{TESSERACT: Eliminating Experimental Bias in Malware Classification across Space and Time.}
USENIX Security, 2019.

\bibitem{arp2022dos}
D. Arp et al. \textit{Dos and Don'ts of Machine Learning in Computer Security.}
USENIX Security, 2022.

\bibitem{sculley2015debt}
D. Sculley et al. \textit{Hidden Technical Debt in Machine Learning Systems.}
NeurIPS, 2015.

\end{thebibliography}

\end{document}
